\documentclass[12pt,a4paper]{article}

\usepackage[utf8]{inputenc}
\usepackage[T2A]{fontenc}
\usepackage[russian]{babel}

\usepackage{amsmath,amssymb,mathtools}
\usepackage{geometry}
\geometry{margin=2.2cm}
\usepackage{microtype}
\usepackage{parskip}

\newcommand{\Z}{\mathbb{Z}}
\newcommand{\F}{\mathbb{F}}
\DeclareMathOperator{\ord}{ord}
\DeclareMathOperator{\Ker}{Ker}
\DeclareMathOperator{\gcdop}{gcd}
\newcommand{\st}{\,:\,}

\begin{document}

\section*{\centering Лекция 2}

\section{Группы и подгруппы}

\textbf{Определение (группа).}
Множество $G$ с бинарной операцией $G\times G\to G$ называется группой, если:
\begin{enumerate}
  \item (ассоциативность) $(gh)x=g(hx)$ для всех $g,h,x\in G$;
  \item (нейтральный элемент) существует $e\in G$: $ge=eg=g$ для всех $g\in G$;
  \item (обратный элемент) для каждого $g\in G$ существует $g^{-1}\in G$:
  $gg^{-1}=g^{-1}g=e$.
\end{enumerate}

Если $gh=hg$ для всех $g,h\in G$, то $G$ называется \textbf{абелевой} (коммутативной).

\textbf{Примеры.}
\begin{enumerate}
  \item $(\Z,+)$;
  \item $(V,+)$ для векторного пространства $V$;
  \item $(\mathbb{R},+)$;
  \item $(\mathbb{R}^{\ast},\cdot)$;
  \item $(S_n,\circ)$ — группа перестановок;
  \item $GL_n(k)$ — группа обратимых матриц над полем $k$.
\end{enumerate}

\textbf{Определение (подгруппа).}
Подмножество $H\subseteq G$ называется подгруппой ($H\le G$), если:
\begin{enumerate}
  \item $x,y\in H \Rightarrow xy\in H$;
  \item $x\in H \Rightarrow x^{-1}\in H$.
\end{enumerate}
Тогда автоматически $e\in H$.

\textbf{Примеры подгрупп.}
\begin{enumerate}
  \item Для $(\Z,+)$: все подгруппы имеют вид $m\Z$.
  \item Подпространство $U\le V$ является подгруппой в $(V,+)$.
  \item $a\mathbb{R}=\{ax\mid x\in\mathbb{R}\}$ — подгруппа в $(\mathbb{R},+)$.
  \item $\{1\}\le G$ (тривиальная подгруппа).
  \item $A_n\le S_n$ — чётные перестановки.
  \item $\left\{\begin{pmatrix}\lambda&0\\0&\lambda\end{pmatrix}\mid \lambda\ne 0\right\}\le GL_2(k)$.
\end{enumerate}

\subsection{Смежные классы}

Пусть $H\le G$. Определим отношение на $G$:
\[
g_1 \sim g_2 \iff g_1^{-1}g_2\in H.
\]
Это отношение эквивалентности. Класс эквивалентности элемента $g$:
\[
gH=\{gh\mid h\in H\},
\]
называется \textbf{левым смежным классом}.

\section{Теорема Лагранжа и следствие про $\varphi(n)$}

\textbf{Разбиение на смежные классы.}
Существуют представители $g_i$ такие, что
\[
G=\bigsqcup_{i\in I} g_iH.
\]

\textbf{Теорема Лагранжа.}
Если $|G|<\infty$ и $H\le G$, то
\[
|H|\mid |G|.
\]
Эквивалентно: $|G|=|H|\cdot [G:H]$, где $[G:H]$ — число левых смежных классов.

\textbf{Замечание.}
Для группы $(\Z/n\Z)^{\ast}$ имеем $\varphi(n)=|(\Z/n\Z)^{\ast}|$.
Из теоремы Лагранжа следует: для $a\in(\Z/n\Z)^{\ast}$ выполнено
\[
\ord(a)\mid \varphi(n)\quad\Rightarrow\quad a^{\varphi(n)}\equiv 1\pmod n.
\]

\section{Гомоморфизмы и изоморфизмы}

\textbf{Определение (гомоморфизм).}
Отображение $f:G\to H$ называется гомоморфизмом, если
\[
f(g_1g_2)=f(g_1)f(g_2)\quad \forall g_1,g_2\in G.
\]

\textbf{Следствия.}
\[
f(e_G)=e_H,\qquad f(x^{-1})=f(x)^{-1}.
\]

\textbf{Определения.}
\[
\Ker(f)=\{g\in G\mid f(g)=e_H\},\qquad
\Im(f)=\{f(g)\mid g\in G\}.
\]
Если $f$ биективно, то это \textbf{изоморфизм}.

\textbf{Примеры.}
\begin{enumerate}
  \item $\Z \to \Z/n\Z$, \; $x\mapsto \overline{x}$.
  \item Для линейного отображения $A:V_1\to V_2$ имеем индуцированное отображение
  $\Lambda^{\dim V_1}(V_1)\to \Lambda^{\dim V_1}(V_2)$ и определитель $\det A$.
  \item $\det:GL_n(k)\to k^{\ast}$ — гомоморфизм.
\end{enumerate}

\textbf{Наблюдение.}
Для $h\in \Im(f)$ множество $f^{-1}(h)$ либо пусто, либо является смежным классом по $\Ker(f)$.
В частности, $f(g\cdot \Ker(f))=f(g)$.

\textbf{Фактор-группа.}
Множество смежных классов $G/\Ker(f)$ с операцией
\[
(g_1\Ker(f))\cdot(g_2\Ker(f))=(g_1g_2)\Ker(f)
\]
образует группу.

\textbf{Теорема об изоморфизме.}
\[
G/\Ker(f)\cong \Im(f).
\]

\section{Цикличность $(\Z/n\Z)^{\ast}$ и первообразные корни}

\textbf{Определение.}
Элемент $g\in(\Z/n\Z)^{\ast}$ называется \textbf{первообразным корнем по модулю $n$},
если
\[
\ord(g)=\varphi(n).
\]

\textbf{Примеры.}
\[
(\Z/7\Z)^{\ast}=\{1,2,3,4,5,6\},\quad \varphi(7)=6,
\]
первообразные корни: $3,5$.

\[
(\Z/8\Z)^{\ast}=\{1,3,5,7\}\cong \Z_2\times \Z_2,
\]
поэтому группа не циклическая.

\textbf{Утверждение.}
Если $g$ — первообразный корень по модулю $n$, то отображение
\[
f:\ (\Z/\varphi(n)\Z,+)\to (\Z/n\Z)^{\ast},\qquad f(m)=g^m
\]
— изоморфизм групп.

\subsection{Леммы для случая простого модуля}

Пусть $p$ — простое.

\textbf{Лемма.}
Если $d\mid (p-1)$, то сравнение $x^d\equiv 1\pmod p$ имеет ровно $d$ решений.

\textbf{Лемма.}
Если $q$ — простое и $q^s\mid (p-1)$, то в $(\Z/p\Z)^{\ast}$ существует элемент порядка $q^s$.

\textbf{Следствие.}
Группа $(\Z/p\Z)^{\ast}$ циклическая, то есть существует первообразный корень по модулю $p$.

\section{Пример КТО}

Найти $x$, если
\[
\begin{cases}
x\equiv -7 \pmod{11}\\
x\equiv -5 \pmod{13}\\
x\equiv 8 \pmod{14}
\end{cases}
\]
Пусть $m_1=11$, $m_2=13$, $m_3=14$,
\[
M=11\cdot 13\cdot 14=2002,\quad
M_1=\frac{M}{11}=182,\quad
M_2=\frac{M}{13}=154,\quad
M_3=\frac{M}{14}=143.
\]
Найдём $\alpha_i$ из $\alpha_i M_i\equiv 1\pmod{m_i}$:
\[
182\equiv 6\pmod{11}\Rightarrow 6\alpha_1\equiv 1\pmod{11}\Rightarrow \alpha_1=2,
\]
\[
154\equiv 11\pmod{13}\Rightarrow 11\alpha_2\equiv 1\pmod{13}\Rightarrow \alpha_2=6,
\]
\[
143\equiv 3\pmod{14}\Rightarrow 3\alpha_3\equiv 1\pmod{14}\Rightarrow \alpha_3=5.
\]
Тогда решение:
\[
x\equiv (-7)\cdot 2\cdot 182 + (-5)\cdot 6\cdot 154 + 8\cdot 5\cdot 143 \pmod{2002}.
\]
Численно:
\[
x\equiv -2548-4620+5720 = -1448 \equiv 554 \pmod{2002}.
\]
Ответ: $x\equiv 554\pmod{2002}$.

\section{Порядок степеней и произведения коммутирующих элементов}

\textbf{Факт.}
Если $\ord(g)=n$, то для любого $s\in\Z$:
\[
\ord(g^s)=\frac{n}{\gcdop(n,s)}.
\]

\textbf{Пример.}
В циклической группе $C_6=\langle g\rangle$:
\[
\ord(g)=6,\quad \ord(g^2)=3,\quad \ord(g^3)=2,\quad \ord(g^4)=3,\quad \ord(g^5)=6.
\]

\textbf{Факт.}
Если $gh=hg$, $\ord(g)=m$, $\ord(h)=n$ и $\gcdop(m,n)=1$, то
\[
\ord(gh)=mn.
\]
Идея: $(gh)^{mn}=g^{mn}h^{mn}=e$; если $(gh)^k=e$, то $g^k=h^{-k}$.
Левая часть имеет порядок, делящий $m$, правая — делящий $n$, при взаимной простоте получается $mn\mid k$.

\end{document}